\begin{figure}[htbp]
	\centering
	\begin{tikzpicture}[>=Stealth, line width=0.8pt]
		
		% ═══════════════════════════════════════════════════════════
		% PANNELLO (a) — Prestazione cinematica dell'incastro
		% ═══════════════════════════════════════════════════════════
		\begin{scope}[xshift=0cm]
			
			%% ── Parametri ──
			\def\xA{0.00}
			\def\yA{0.00}
			\def\Lcamp{0.65}      % era 0.50, semi-altezza fascia campita
			
			%% ── Corpo C (ingrandito) ──
			\draw[line width=1pt] 
			(\xA, -\Lcamp)
			.. controls ( 0.65, -1.65) and ( 1.65, -1.75)
			.. ( 2.25, -1.40)
			.. controls ( 3.00, -0.75) and ( 3.10,  0.50)
			.. ( 2.65, 1.25)
			.. controls ( 1.90, 1.75) and ( 0.75, 1.65)
			.. (\xA, \Lcamp)
			-- cycle;
			\node at (1.80, 0.10) {$\mathcal{C}$};
			
			%% ── Campitura (tratteggio obliquo) lungo il bordo sinistro ──
			\foreach \y in {-0.55,-0.40,-0.25,-0.10,0.05,0.20,0.35,0.50}
			\draw[thin] (\xA, \y) -- ++(-0.13, -0.10);
			
			
			%% ── Punto A e label ──
			\fill (\xA, \yA) circle (1.5pt);
			\node[anchor=south west, font=\small] 
			at ({\xA+0.05}, {\yA+0.05}) {$A$};
			
			%% ── Equazioni di vincolo (a sinistra, fuori dal corpo) ──
			\node[font=\small, anchor=east] 
			at (-0.40, 0.50) {$\ub(A) = \zerovb$};
			\node[font=\small, anchor=east] 
			at (-0.40, 0.10) {$\vartheta = 0$};
			
			%% ── Etichetta pannello ──
			\node[anchor=north, font=\small] 
			at (-0.06, -1.0) {$(a)$};
			
		\end{scope}
		% ═══════════════════════════════════════════════════════════
		% PANNELLO (b) — Prestazione statica dell'incastro
		% ═══════════════════════════════════════════════════════════
		\begin{scope}[xshift=5cm]
			
			%% ── Parametri ──
			\def\xA{0.00}
			\def\yA{0.00}
			\def\Lcamp{0.65}      % era 0.50, semi-altezza fascia campita
			
			%% ── Corpo C (ingrandito) ──
			\draw[line width=1pt] 
			(\xA, -\Lcamp)
			.. controls ( 0.65, -1.65) and ( 1.65, -1.75)
			.. ( 2.25, -1.40)
			.. controls ( 3.00, -0.75) and ( 3.10,  0.50)
			.. ( 2.65, 1.25)
			.. controls ( 1.90, 1.75) and ( 0.75, 1.65)
			.. (\xA, \Lcamp)
			-- cycle;
			\node at (1.80, 0.10) {$\mathcal{C}$};
			
			%% ── Campitura (tratteggio obliquo) lungo il bordo sinistro ──
			\foreach \y in {-0.55,-0.40,-0.25,-0.10,0.05,0.20,0.35,0.50}
			\draw[thin] (\xA, \y) -- ++(-0.13, -0.10);
			
			%% ── Punto A e label ──
			\fill (\xA, \yA) circle (1.5pt);
			\node[anchor=south east, font=\small] 
			at ({\xA-0.05}, {\yA+0.05}) {$A$};
			
			%% ──────────────────────────────────────────────────
			%% REAZIONI: X(A), Y(A), μ(A)
			%% ──────────────────────────────────────────────────
			
			%% Forza X(A): da A verso destra
			\draw[->, line width=1.4pt] 
			(\xA+0.05, \yA) -- (\xA+1.05, \yA);
			\node[anchor=north, font=\small] at ({\xA+0.85}, {\yA-0.05}) {$X(A)\,\ab_x$};
			
			%% Forza Y(A): da A verso l'alto
			\draw[->, line width=1.4pt] 
			(\xA, \yA+0.05) -- (\xA, \yA+1.05);
			\node[anchor=south, font=\small] 	at ({\xA-0.05}, {\yA+1.0}) {$Y(A)\,\ab_y$};
			
			%% ── Terna x, y in basso a sinistra ──
			\draw[->, line width=0.8pt] (-1.50, -1.20) -- (-0.70, -1.20);
			\node[anchor=west, font=\footnotesize] at (-0.70, -1.25) {$x$};
			\draw[->, line width=0.8pt] (-1.50, -1.20) -- (-1.50, -0.40);
			\node[anchor=south, font=\footnotesize] at (-1.50, -0.40) {$y$};
			
			%% Coppia μ(A): arco antiorario dentro al corpo
			\def\xMu{1.10}     % era 0.80
			\def\yMu{0.50}     % era 0.80
			
			%% Coppia μ(A): arco di π/2 centrato in A, fra le due frecce di forza
			\def\rmu{0.55}
			\draw[->, line width=1.4pt] 
			({\xA+\rmu}, \yA) arc(0:90:\rmu);
			\node[font=\small] 
			at ({\xA+\rmu*1.7}, {\yA+\rmu*0.707+0.05}) {$\mu(A)$};
			
			
			%% ── Etichetta pannello ──
			\node[anchor=north, font=\small] 	at (-0.06, -1.0) {$(b)$};
			
		\end{scope}
	\end{tikzpicture}
	\caption[Incastro in $A$: prestazioni cinematica e statica.]{%
		Incastro applicato al corpo $\mathcal{C}$ nel punto $A$.
		Prestazione cinematica (a): nessuno spostamento è ammesso, 
		né traslazione né rotazione ($\ub(A) = \zerovb$, 
		$\vartheta = 0$); il corpo è completamente bloccato.
		Prestazione statica (b): la reazione vincolare è una forza
		generica $X(A)\,\ab_x + Y(A)\,\ab_y$ e una coppia $\mu(A)$,
		applicate in $A$, in accordo con la molteplicità tre.}
	\label{fig:incastro}
\end{figure}
